\documentclass[11pt,a4paper]{article}

\usepackage[utf8]{inputenc}
\usepackage[T1]{fontenc}
\usepackage[spanish]{babel}
\usepackage[left=1.5cm, right=1.5cm, top=1.5cm, bottom=1.5cm]{geometry}
\usepackage{hyperref}
\usepackage{fontawesome5}
\usepackage{enumitem}
\usepackage{titlesec}
\usepackage{parskip}
\usepackage{xcolor}
\usepackage{array}
\usepackage{tabularx}

\definecolor{sectioncolor}{RGB}{30, 100, 180}

\hypersetup{
    colorlinks=true,
    urlcolor=sectioncolor,
    linkcolor=sectioncolor
}

\titleformat{\section}
  {\large\bfseries\color{sectioncolor}}
  {}{0em}{}[\titlerule]

\setlength{\parindent}{0pt}
\setlength{\parskip}{4pt}

\begin{document}

\pagestyle{empty}

% ─── Encabezado ────────────────────────────────────────────────────────────────
\begin{center}
  {\LARGE\textbf{Pablo Fernando Badilla Torrealba}} \\[4pt]
  \faEnvelope\ \href{mailto:pbadilla.torrealba@gmail.com}{pbadilla.torrealba@gmail.com}
  \quad
  \faLinkedin\ \href{https://www.linkedin.com/in/pablo-badilla-torrealba-473b2315b/}{pablo-badilla-torrealba}
  \quad
  \faPhone\ +56 9 7139 6025
  \quad
  \faGithub\ \href{https://github.com/pbadillatorrealba}{pbadillatorrealba}
\end{center}

% ─── Resumen Profesional ────────────────────────────────────────────────────────
\section{Resumen Profesional}

Ingeniero Civil en Computación y Magíster en Ciencias de la Computación especializado en
Ciencia de Datos, Machine Learning e IA Generativa, con más de 7 años de experiencia
liderando e implementando sistemas productivos para banca, retail, academia y start-ups,
entre otros.

Ha liderado equipos técnicos con enfoque en gestión de clientes, arquitectura de sistemas
escalables y mentoría, promoviendo la creación de productos de alto estándar, buenas
prácticas y ejecución sólida de principio a fin.

% ─── Antecedentes Laborales ────────────────────────────────────────────────────
\section{Antecedentes Laborales}

\textbf{Iniciativa de Datos e Inteligencia Artificial (IDIA) -- Universidad de Chile} \hfill Santiago, Chile \\
\textit{Tech Lead, Área de Transferencia Tecnológica} \hfill Abril 2026 -- Presente
\begin{itemize}[leftmargin=1.5em, itemsep=1pt, topsep=3pt]
  \item \textbf{Construcción y Definición del Área:} Liderazgo en la definición y
        estructuración del área de transferencia tecnológica dentro de IDIA, iniciativa que
        reúne a los principales actores académicos de la FCFM (DIE, DII, DCC, DIM) para
        articular capacidades en investigación, transferencia, docencia y proyectos.
  \item \textbf{ETL basado en IA -- Ministerio de las Culturas:} Desarrollo de sistema de
        extracción y transformación automática de datos de indicadores (ETL basado en IA)
        usando Python, LLMs/IA Generativa y FastAPI, permitiendo al Ministerio procesar
        automáticamente datos de indicadores culturales.
\end{itemize}

\vspace{6pt}
\textbf{Banco Falabella} \hfill Santiago, Chile \\
\textit{Data Scientist} \hfill Agosto 2024 -- Abril 2026
\begin{itemize}[leftmargin=1.5em, itemsep=1pt, topsep=3pt]
  \item \textbf{Productos y Beneficios:} Chat con Agentes de IA para consultas sobre
        productos, beneficios, operaciones y sucursales en Chile, Colombia, México y Perú.
        +100k atenciones mensuales, reducción del 80\% de costos, latencia y experiencia de
        usuario mejoradas. Gestión de relaciones con clientes internos en todos los países en
        colaboración con equipo de Canales Digitales.
  \item \textbf{Desconocimiento de transacciones:} Desarrollo y puesta en marcha del
        proyecto enfocado en iniciar un proceso de desconocimiento de transacción por fraude
        en transacciones no conciliadas. Primera iniciativa de IA publicada por el banco,
        abriendo las puertas a proyectos más complejos y a estrategia basada en IA.
        Colaboración con equipo de Canales Digitales.
  \item \textbf{Librería Servidor Concurrente de Conversaciones:} Diseño e implementación
        de librería basada en algoritmos concurrentes y distribuidos para coordinar
        conversaciones en distintos casos de uso de IA. Base para los nuevos proyectos
        basados en IA conversacional.
  \item \textbf{Mentoría:} Mentoría en desarrollo conceptual de agentes conversacionales
        basados en IA, tecnologías (LangChain/LangGraph/LangFuse, RAG, etc.), tecnologías
        Backend (FastAPI, SQLAlchemy, Pydantic, Asyncio, etc.) y uso de IA aplicada al
        desarrollo (Copilot) al equipo.
\end{itemize}

\vspace{6pt}
\textbf{Walmart Chile} \hfill Santiago, Chile \\
\textit{Machine Learning Engineer} \hfill Abril 2023 -- Abril 2024
\begin{itemize}[leftmargin=1.5em, itemsep=1pt, topsep=3pt]
  \item \textbf{Call Center Timeseries:} Desarrollo y despliegue de modelo de timeseries
        basado en Machine Learning para la predicción de demanda diaria del Call Center.
        Reducción del error de un 60\% a un 20\%. Evaluación mediante A/B testing y shadow
        deployment con métricas de ML y KPIs asociados a costos. Planificación de hitos y
        presentación conjunta de avances y resultados a clientes internos del equipo de
        operaciones junto con PO.
  \item \textbf{Simulación de tiempos de espera en Colas:} Optimización de proyecto para
        medir tiempos de espera en colas según formato (Lider, Express, Acuenta, Central
        Mayorista) y tipo de caja (MM1 y MMS) usando teoría de colas. Refactorización
        sustancial de arquitectura y rendimiento del sistema.
  \item \textbf{Customer Lifetime Value (CLV):} Reimplementación del modelo Customer
        Lifetime Value para el área de Marketing, permitiendo campañas personalizadas según
        segmento de cliente detectado.
\end{itemize}

\vspace{6pt}
\textbf{DashAI} \hfill Santiago, Chile \\
\textit{Tech Lead} \hfill Enero 2023 -- Junio 2024
\begin{itemize}[leftmargin=1.5em, itemsep=1pt, topsep=3pt]
  \item \textbf{Liderazgo del Equipo de Desarrollo:} Liderazgo técnico de equipo de 5
        desarrolladores (memoristas Universidad de Chile). Responsable de guiar diversos
        proyectos e iniciativas en aspectos técnicos y arquitectónicos.
  \item \textbf{Gestión y Coordinación:} Dirección del equipo, ejecución de ceremonias
        (dailies, plannings, refinamientos, etc.), definición de requisitos técnicos y de
        integración. Repositorio:
        \href{https://github.com/DashAISoftware/DashAI}{github.com/DashAISoftware/DashAI}
  \item \textbf{Postulación a fondos y Transferencia Tecnológica:} Postulación a FONDEF
        2023. Generación de modelo de negocios y plan para productivizar la aplicación y el
        framework conceptual de la plataforma. Diseño de estrategia de transferencia
        tecnológica, análisis de competidores y modelo de negocios.
  \item \textbf{Construcción de componentes Core:} Creación de Registro de Componentes
        (Plugins), inyección de dependencias, integración de subproyectos, test masivos,
        Front-end inicial, buenas prácticas de desarrollo (CI, tests, verificación estática
        de código, tipado), etc.
\end{itemize}

\vspace{6pt}
\textbf{BCI} \hfill Santiago, Chile \\
\textit{Data Scientist} \hfill Febrero 2023 -- Mayo 2023
\begin{itemize}[leftmargin=1.5em, itemsep=1pt, topsep=3pt]
  \item \textbf{Features para evaluación de riesgo crediticio:} Creación de features
        basadas en comportamiento transaccional de clientes para explorar mejoras en modelos
        de evaluación de riesgo crediticio.
\end{itemize}

\vspace{6pt}
\textbf{Centro Nacional de Inteligencia Artificial -- CENIA} \hfill Santiago, Chile \\
\textit{Machine Learning Engineer} \hfill Enero 2022 -- Enero 2023
\begin{itemize}[leftmargin=1.5em, itemsep=1pt, topsep=3pt]
  \item \textbf{Modelo ML End-to-End para VRP:} Diseño y despliegue de modelo de Machine
        Learning que permite incorporar preferencias de usuario en el cálculo de soluciones
        del Problema de Enrutamiento de Vehículos (VRP). Comunicación con cliente,
        presentaciones de avances, gestión de hitos y entregas en colaboración con equipo de
        Optimización. Entrega de paquete completo listo para despliegue en sistemas internos.
\end{itemize}

\vspace{6pt}
\textbf{Universidad de Chile} \hfill Santiago, Chile \\
\textit{Profesor Part-Time} \hfill Marzo 2021 -- Julio 2024
\begin{itemize}[leftmargin=1.5em, itemsep=1pt, topsep=3pt]
  \item \textbf{Docente Laboratorio de Programación Científica para Ciencia de Datos:}
        Docente del curso de Laboratorio de Programación Científica para Ciencia de Datos
        para el Magíster en Ciencia de Datos (MDS) de la Universidad de Chile.
  \item \textbf{Docente Manejo y Extracción de Datos en Python:} Docente del curso de
        Manejo y Extracción de Datos en Python del Diploma en Python aplicado a la Ciencia
        de Datos del Departamento de Ciencias de la Computación de la Universidad de Chile.
\end{itemize}

\vspace{6pt}
\textbf{Kwiyx SPA} \hfill Santiago, Chile \\
\textit{Desarrollador Full Stack} \hfill Abril 2020 -- Enero 2021
\begin{itemize}[leftmargin=1.5em, itemsep=1pt, topsep=3pt]
  \item \textbf{App de inversión inmobiliaria:} Creación del front-end de aplicación de
        inversión inmobiliaria.
  \item \textbf{Gestión de recursos humanos en hospitales:} Desarrollo full-stack de
        aplicación para la gestión de recursos humanos en hospitales en pandemia basada en
        optimización de turnos.
\end{itemize}

\vspace{6pt}
\textbf{Desarrollador Freelance} \hfill Santiago, Chile \\
\textit{Data Scientist y Desarrollador Full Stack} \hfill Julio 2019 -- Enero 2021
\begin{itemize}[leftmargin=1.5em, itemsep=1pt, topsep=3pt]
  \item \textbf{Sistema de Detección de Tópicos -- Universidad de Chile \& Fundación Tribu:}
        Construcción de sistema de detección de tópicos usando embeddings de Transformers
        para análisis del encuentro participativo 2019 de ``Lxs 400'' en torno a la nueva
        Constitución.
  \item \textbf{Detección Automática de Texto -- Sotta y CIA:} Desarrollo de sistema basado
        en OCR + Regex para la extracción automática de texto de pagarés de deudas de
        créditos automotrices.
\end{itemize}

\vspace{6pt}
\textbf{Acción Circular} \hfill Santiago, Chile \\
\textit{Desarrollador Full Stack} \hfill Abril 2019 -- Julio 2019
\begin{itemize}[leftmargin=1.5em, itemsep=1pt, topsep=3pt]
  \item \textbf{Sistema de gestión de residuos reciclables:} Diseño e implementación de
        software para la gestión de recolección de residuos reciclables.
\end{itemize}

\vspace{6pt}
\textbf{Cognus Consulting} \hfill Santiago, Chile \\
\textit{Data Analyst Intern} \hfill Diciembre 2017 -- Febrero 2018
\begin{itemize}[leftmargin=1.5em, itemsep=1pt, topsep=3pt]
  \item \textbf{Solución de Data Warehousing y Business Intelligence:} Construcción de
        solución de Data Warehousing, Business Intelligence y visualizaciones usando Pentaho.
\end{itemize}

% ─── Proyectos ──────────────────────────────────────────────────────────────────
\section{Proyectos}

\textbf{Word Embeddings Fairness Evaluation Framework} (Enero 2020 -- Junio 2025): \\
Liderazgo y desarrollo de una librería para la medición y mitigación de sesgos de género,
étnicos y religiosos en modelos de embeddings de palabras (precursores a modelos de
lenguaje). 2 papers publicados, 182 estrellas en GitHub. \\
Repositorio: \href{https://github.com/dccuchile/wefe}{github.com/dccuchile/wefe} \quad
Documentación: \href{https://wefe.readthedocs.io}{wefe.readthedocs.io}

% ─── Antecedentes Académicos ───────────────────────────────────────────────────
\section{Antecedentes Académicos}

\textbf{Magíster en Ciencias de la Computación} \hfill Santiago, Chile \\
\textit{Universidad de Chile} \hfill 2019 -- 2020

\vspace{4pt}
\textbf{Ingeniería Civil en Computación} \hfill Santiago, Chile \\
\textit{Universidad de Chile} \hfill 2013 -- 2019

% ─── Publicaciones ─────────────────────────────────────────────────────────────
\section{Publicaciones}

\textbf{WEFE: A Python Library for Measuring and Mitigating Bias in Word Embeddings} \\
Pablo Badilla, Felipe Bravo-Marquez, María José Zambrano, Jorge Pérez. \\
\textit{Journal of Machine Learning Research}, Vol. 26, No. 156, pp. 1--6, 2025. \\
\href{http://jmlr.org/papers/v26/22-1133.html}{http://jmlr.org/papers/v26/22-1133.html}

\vspace{6pt}
\textbf{WEFE: The Word Embeddings Fairness Evaluation Framework} \\
Pablo Badilla, Felipe Bravo-Marquez, Jorge Pérez. \\
\textit{Proceedings of the Twenty-Ninth International Joint Conference on Artificial
Intelligence (IJCAI-20)}, pp. 430--436, 2020. \\
\href{https://doi.org/10.24963/ijcai.2020/60}{https://doi.org/10.24963/ijcai.2020/60}

% ─── Habilidades y Herramientas ────────────────────────────────────────────────
\section{Habilidades y Herramientas}

\begin{itemize}[leftmargin=1.5em, itemsep=2pt, topsep=3pt]
  \item \textbf{Idiomas:} Español Nativo, Inglés B2 TOEFL-ITP
  \item \textbf{Lenguajes de Programación:} Python (Avanzado), Javascript (Intermedio)
  \item \textbf{IA Generativa:} LangChain, LangGraph, LangFuse, APIs de OpenAI, Gemini, Azure
  \item \textbf{Ciencia de Datos:} Numpy, Pandas, SQL, BigQuery, Spark, JupyterLab, Plotly, DataBricks
  \item \textbf{Machine Learning:} Scikit-Learn, XGBoost, LightGBM, PyTorch, HuggingFace
  \item \textbf{Desarrollo Backend:} FastAPI, SQLAlchemy, Pydantic, Asyncio, Concurrencia, Inyección de Dependencias
  \item \textbf{Cloud e Infraestructura:} Google Cloud Platform, VertexAI, Docker, Kubeflow
  \item \textbf{Herramientas y Metodologías:} Kedro, Github Actions, Gitlab CI, Github Copilot, MLOps, Scrum
  \item \textbf{Sistemas Operativos y Ofimática:} Windows, Linux, macOS, Microsoft Office, Google Docs, Sharepoint, OneDrive, Google Drive
\end{itemize}

% ─── Información Adicional ─────────────────────────────────────────────────────
\section{Información Adicional}

\begin{itemize}[leftmargin=1.5em, itemsep=2pt, topsep=3pt]
  \item Graduado de Ingeniería y Magíster con Distinción Máxima.
  \item Estudiante destacado años 2016 y 2018: Facultad de Ciencias Físicas y Matemáticas,
        Universidad de Chile.
  \item Guía de Memoria de Título: Implementación de módulo de mitigación de sesgo en
        modelos de embeddings de palabras en WEFE.
  \item Docente Invitado: Charlas sobre Machine Learning Operations para varios cursos de
        pre y postgrado en la Universidad de Chile.
  \item \textbf{Microsoft College Code Competition:} Tercer lugar en Microsoft College Code
        Competition 2018.
  \item \textbf{Revisor de Papers:} Revisor secundario de la conferencia: The Joint
        Conference of the 59th Annual Meeting of the Association for Computational Linguistics
        and the 11th International Joint Conference on Natural Language Processing
        (ACL-IJCNLP 2021).
\end{itemize}

\end{document}
